% from ega4-2-fr.tex
\oldpage[IV-2]{37}
\item[\emph{b)}] Il existe un voisinage ouvert $U$ de $x$ et une section
$f\in\Gamma(U,\mathcal F)$ tels que $U\cap Z$ soit égal à
$\operatorname{Supp}(\mathcal O_U\cdot f)$.
\item[\emph{b$'$)}] Il existe un voisinage ouvert $U$ de $x$ et une section
$f\in\Gamma(U,\mathcal F)$ tels que $U\cap Z$ soit une composante
irréductible de $\operatorname{Supp}(\mathcal O_U\cdot f)$.
\item[\emph{c)}] Il existe un voisinage ouvert $U$ de $x$ et un sous-Module
de $\mathcal F|U$ isomorphe à $\mathcal O_Z|U$ ($\mathcal O_Z$ étant
identifié à un quotient de $\mathcal O_X$).
\item[\emph{c$'$)}] Il existe un voisinage ouvert $U$ de $x$ et un
sous-Module cohérent $\mathcal G$ de $\mathcal F|U$ tels que $U\cap Z$ soit
une composante irréductible de $\operatorname{Supp}(\mathcal G)$.
\end{enumerate}
\end{proposition}

Il est clair que \emph{c)} implique \emph{b)}, car il suffit de prendre pour
$f$ l'élément de $\Gamma(U,\mathcal F)$ qui correspond à la section unité de
$\mathcal O_Z|U$. Comme $U\cap Z$ est irréductible $(0_{\mathrm I},2.1.6)$,
\emph{b)} implique \emph{b$'$)} et \emph{b$'$)} implique \emph{c$'$)} puisque
$\mathcal O_X$ est cohérent $(0_{\mathrm I},5.3.4)$. Pour voir que
\emph{c$'$)} implique \emph{a)}, on peut se borner au cas où
$U=X=\operatorname{Spec}(A)$ est affine, $A$ étant donc noethérien, et où
$\mathcal F=\widetilde M$, où $M$ est un $A$-module, et
$\mathcal G=\widetilde N$, où $N$ est un sous-module de $M$, de type fini.
On sait alors que les éléments minimaux de $\operatorname{Supp}(\mathcal G)$
sont les points maximaux de $V(\operatorname{Ann}(N))$ $(0_{\mathrm I},1.7.4)$,
et ce sont aussi les éléments minimaux de $\operatorname{Ass}(N)$ (Bourbaki,
\emph{Alg. comm.}, chap.~IV, \S~1, n\textsuperscript{o}~3, cor.~1 de la
prop.~7)~; comme $\operatorname{Ass}(N)\subset\operatorname{Ass}(M)
=\operatorname{Ass}(\mathcal F)$, on voit que \emph{c$'$)} implique
\emph{a)}. Enfin \emph{a)} implique \emph{c)} en vertu de (3.1.2), en
prenant encore $X$ affine, $\mathcal F=\widetilde M$, et $Z$ défini par
l'idéal $fA$ (avec les notations de (3.1.2)).

\begin{corollary}[3.1.4]
\label{IV.2.3.1.4-fr}
Soient $X$ un préschéma localement noethérien, $\mathcal F$ un
$\mathcal O_X$-Module cohérent. Alors les cycles premiers maximaux associés à
$\mathcal F$ sont les composantes irréductibles de
$\operatorname{Supp}(\mathcal F)$, et les points génériques de ces composantes
sont les points $x\in X$ tels que $\mathcal F_x$ soit un
$\mathcal O_x$-Module de longueur finie et $\ne0$.
\end{corollary}

En effet, si $x$ est le point générique d'une des composantes irréductibles
$Z$ de $\operatorname{Supp}(\mathcal F)$, il résulte de l'équivalence de
\emph{a)} et \emph{c$'$)} dans (3.1.3) que $x$ appartient à
$\operatorname{Ass}(\mathcal F)$, et $Z$ est un cycle premier associé à
$\mathcal F$, nécessairement maximal en vertu de (3.1.1.1)~; la réciproque
résulte trivialement de (3.1.1.1)~; enfin, la dernière assertion étant
évidemment locale résulte de Bourbaki, \emph{Alg. comm.}, chap.~IV, \S~2,
n\textsuperscript{o}~5, cor.~2 de la prop.~7.

\begin{corollary}[3.1.5]
\EGAReaderTarget{ega4.c04.cor.3-1-5}{IV.2.3.1.5-fr}
Soient $X$ un préschéma localement noethérien, $\mathcal F$ un
$\mathcal O_X$-Module quasi-cohérent. Pour que $\mathcal F=0$, il faut et il
suffit que $\operatorname{Ass}(\mathcal F)=\varnothing$.
\end{corollary}

La question étant locale, on est ramené au cas où $X$ est affine, et la
conclusion résulte aussitôt de (3.1.2) et de Bourbaki, \emph{Alg. comm.},
chap.~IV, \S~1, n\textsuperscript{o}~1, cor.~1 de la prop.~2.

\begin{proposition}[3.1.6]
\label{IV.2.3.1.6-fr}
Soient $X$ un préschéma localement noethérien, $\mathcal F$ un
$\mathcal O_X$-Module cohérent~; alors $\operatorname{Ass}(\mathcal F)$ est
\emph{localement fini} (c'est-à-dire que tout point de $X$ admet un voisinage
dont l'intersection avec $\operatorname{Ass}(\mathcal F)$ est finie).
\end{proposition}

Il suffit de considérer le cas où $X$ est affine, donc noethérien, et alors la
proposition résulte de (3.1.2) et de Bourbaki, \emph{Alg. comm.}, chap.~IV,
\S~1, n\textsuperscript{o}~4, cor. du th.~2.

